% Part: first-order-logic
% Chapter: sequent-calculus
% Section: rules-and-proofs

\documentclass[../../../include/open-logic-section]{subfiles}

\begin{document}

\olfileid[hi]{mvl}{seq}{rul}

\olsection{नियम और \usetoken{P}{derivation}}

आगे $\Gamma, \Delta, \Pi, \Lambda$ को !!{sentence}ों के परिमित अनुक्रम मानें।

\begin{defn}[सीक्वेंट]
\emph{$n$-पक्षीय सीक्वेंट} इस रूप का व्यंजक है:
\[
\Gamma_1 \nSequent \dots \nSequent \Gamma_n,
\]
जहाँ प्रत्येक $\Gamma_1$, भाषा~$\Lang L$ के !!{sentence}ों का कोई परिमित
(संभवतः रिक्त) अनुक्रम है।
\end{defn}

\begin{defn}[आरंभिक सीक्वेंट]
किसी भी !!{sentence}~$!A$ के लिए $!A \nSequent \dots \nSequent !A$ रूप का
$n$-पक्षीय सीक्वेंट \emph{$n$-पक्षीय आरंभिक सीक्वेंट} कहलाता है।

यदि भाषा में कोई $0$-स्थानिक संयोजक~$\star$, अर्थात् कोई प्रतिज्ञप्तिक
अचर, हो, तो हम सीक्वेंट $\dots \nSequent \star \nSequent \dots$ को भी
आरंभिक मानते हैं। इसमें $\star$ उस सत्य-मूल्य से जुड़े स्थान में आता है जो
$\tf{\star} \in V$ के संगत है, और अन्य स्थान रिक्त रहते हैं।
\end{defn}

किसी $n$-मूल्य तर्कशास्त्र~$\Log{L}$ के प्रत्येक संयोजक के लिए उस संयोजक
द्वारा $\Log{L}$ में लिए जा सकने वाले प्रत्येक सत्य-मूल्य का एक तार्किक
नियम होता है। $\Log{L}$ के $n$-पक्षीय सीक्वेंट कलन में !!{derivation}याँ
सीक्वेंटों के वृक्ष होती हैं। इनके सबसे ऊपर के सीक्वेंट आरंभिक सीक्वेंट होते
हैं; और यदि कोई सीक्वेंट एक या अधिक अन्य सीक्वेंटों के नीचे आता है, तो उसे
$\Log{L}$ के संयोजकों के किसी अनुमान-नियम से सही ढंग से निकलना चाहिए।

\begin{defn}[प्रमेय]
कोई !!{sentence}~$!A$, $n$-मूल्य तर्कशास्त्र~$\Log{L}$ का \emph{प्रमेय} है,
यदि उस $n$-सीक्वेंट की कोई !!{derivation} हो जिसमें $\Log{L}$ के प्रत्येक
निर्दिष्ट सत्य-मूल्य के संगत स्थान पर $!A$ आता हो। यदि $!A$ प्रमेय है तो हम
$\Proves[\Log{L}] !A$ लिखते हैं, और यदि नहीं है तो
$\Proves/[\Log{L}] !A$।
\end{defn}

\begin{defn}[!!^{derivability}]
कोई !!{sentence}~$!A$, $n$-मूल्य तर्कशास्त्र~$\Log{L}$ में !!{sentence}ों के
समुच्चय~$\Gamma$ से \emph{!!{derivable}} है,
$\Gamma \Proves[\Log{L}] !A$, तभी जब कोई परिमित उपसमुच्चय
$\Gamma_0 \subseteq \Gamma$ और $\Gamma_0$ के !!{sentence}ों का कोई अनुक्रम
$\Gamma_0'$ ऐसा हो कि निम्न सीक्वेंट की कोई !!{derivation} हो:
\[
\Lambda_1 \nSequent \dots \nSequent \Lambda_n,
\]
जहाँ स्थान $i$ किसी निर्दिष्ट सत्य-मूल्य के संगत हो तो $\Lambda_i$, $!A$
है, और अन्यथा $\Gamma_0'$ है। यदि $!A$, $\Gamma$ से !!{derivable} नहीं
है, तो हम $\Gamma \Proves/ !A$ लिखते हैं।
\end{defn}

उदाहरणार्थ, त्रिमूल्य लुकासिएविच तर्कशास्त्र का $3$-पक्षीय सीक्वेंट कलन
होता है। किसी $3$-पक्षीय सीक्वेंट
$\Gamma \nSequent \Pi \nSequent \Delta$ में $\Gamma$, $\False$ के;
$\Delta$, $\True$ के; और $\Pi$, $\Undef$ के संगत होता है। अभिगृहीत
$!A \nSequent !A \nSequent !A$ हैं। चूँकि केवल $\True$ निर्दिष्ट है,
$\Gamma \Proves[\LogLuk[3]] !A$ तभी जब सीक्वेंट
$\Gamma \nSequent \Gamma \nSequent !A$ की कोई !!{derivation} हो।
(यदि $\Undef$ भी निर्दिष्ट होता, तो हमें
$\Gamma \nSequent !A \nSequent !A$ की !!{derivation} चाहिए होती।)

\end{document}
